\section{Limitation and Future Work}
\label{sec:limitation_and_future_work}




In this work, we introduce \framework that facilitates multiple autonomous agents to simulate human groups to accomplish tasks, and discuss the emergent social behaviors of agents during this process. \framework is an advanced attempt; thus, there are some techniques within \framework that still have room for improvement and are worthy of exploration. In this section, we delve into these aspects for further illustration.


\textbf{More Capable Agents and More Challenging Scenarios.}
The \framework is designed to enable various multiple LLM-based agents to collaboratively accomplish tasks. In the current research, we have utilized state-of-the-art agents based on GPT-4. With the advancements in LLMs, such as the newly released version of ChatGPT that incorporates voice and image capabilities \citep{OpenAI_ChatGPT}, LLM-based agents have more perceptual capabilities, including seeing, hearing, and speaking. These enhancements may increase the potential of agents and allow them to accomplish more complex real-world tasks based on the \framework framework.

\textbf{Multi-party Communication Among Agents.}
The currently proposed autonomous agents \citep{richards2023auto, nakajima2023babyagi, agentgpt, DBLP:journals/corr/abs-2305-16291} LLMs possess excellent instruction comprehension capabilities \citep{wei2022finetuned,NEURIPS2020_1f89885d}. This enables them to follow given human instructions and accomplish tasks within a one-on-one (human-to-AI) scenario. However, multi-agent collaboration involves a \textit{multi-party communication} \citep{wei2023multiparty} scenario that requires the capability to autonomously determine \textit{when to speak} and \textit{whom to speak}. This leads to difficulties in communication among the agents during the collaborative decision-making step within the \framework framework. Hence, there are two directions worth exploring. Firstly, akin to the aforementioned, we can explore more effective mechanisms for managing agent communication. Additionally, we can design more advanced perceptual-aware LLMs \citep{OpenAI_ChatGPT} that can autonomously interact with their environments\footnote{This kind of perceptual-aware agent has long been a goal of embodied AI \citep{ahn2022i,DBLP:conf/icml/DriessXSLCIWTVY23}, which is a promising direction to explore.}, including other agents.


\textbf{Leverage Emergent Behaviors and Mitigate Safety Issues.}
In Section \ref{ssec:emerging_behaviors}, we identified both emergent positive and harmful behaviors. Exploring ways to leverage positive behaviors for improving work efficiency and effectiveness, as well as mitigating harmful behaviors, are promising directions.
